\documentclass[10pt,a4paper]{amsart}
\usepackage[margin=19mm]{geometry}
\usepackage{amsmath,amssymb,mathtools}
\usepackage[scr=rsfs,cal=euler]{mathalfa}
\usepackage{xfrac}
\usepackage[unicode,hidelinks,bookmarks=false]{hyperref}
\newcommand{\ie}{i.e.\ }
\newcommand{\cf}{cf.\ }
\newcommand{\resp}{respectively\ }
\newcommand{\Subsection}{\subsection}
\providecommand{\nobreakdash}{\nobreak\hbox{-}}
\setlength{\emergencystretch}{2em}
\multlinegap0pt
\newtheorem{comment}{Comment}
\usepackage{amssymb}
\usepackage{cancel}
\newtheorem{lemma}{Lemma}
\newtheorem{claim}{Claim}
\newcommand{\CH}{\mathsf{CH}}
\newcommand{\forceP}{\mathbb{P}}
\newcommand{\forceQ}{\mathbb{Q}}
\title{A Failure of $\Pi^1_{n+3}$-Reduction in the Presence of $\Sigma^1_{n+3}$-Separation}
\author{Stefan Hoffelner}
\date{Source: arXiv:2312.02540v2 (CC BY 4.0)}
\begin{document}
\maketitle

\noindent\emph{Source and licence.} A Failure of $\Pi^1_{n+3}$-Reduction in the Presence of $\Sigma^1_{n+3}$-Separation. Version arXiv:2312.02540v2 (CC BY 4.0), distributed by arXiv under the Creative Commons Attribution 4.0 International licence, http://creativecommons.org/licenses/by/4.0/, as recorded in the arXiv licence metadata for this exact version and shown on the abstract page https://arxiv.org/abs/2312.02540v2. Full licence notice: \texttt{licenses/arxiv-2312.02540-CC-BY-4.0.txt}.\par\smallskip

\noindent\textit{Editorial scope.} Counted unit: the article's lemma productallowable (source0.tex lines 650--655). The article's complete original proof of it is delivered with it (source0.tex lines 656--811, one \texttt{proof} environment of 156 lines). No mathematics is written, repaired or re-derived: the statement and the proof are the article's own lines, in the article's own notation, and no retained line is substituted or re-flowed.  Retained uncounted as the source-local context the proof needs: lines 416--425.  Nothing was omitted from any retained span, so the package carries no omission record.  No retained line contains a citation, so the wrapper carries no bibliography and no bibliography entry.  No retained line contains a figure environment, an image-inclusion command or drawing code, so no image is delivered and the package declares no asset dependency.  Editorial formatting only, in addition: an AMS \texttt{amsart} preamble that restates the article's own declarations and macros used by the retained lines, namely the environments \texttt{lemma}, \texttt{claim} on separate counters, together with the standard packages those lines need.  The retained environments print in this document as Lemma 1, Claim 1 in order of appearance, whereas the article prints the same items with the numbering of its own full document.  The complete native source of the article as distributed in the arXiv e-print for this exact version is bundled unchanged as \texttt{sources/151021/source0.tex}.
 \begin{lemma} \label{FirstPropertiesOfAllowableForcings}


\begin{enumerate}
\item If $\forceP=(\forceP(\beta) \, : \, \beta < \delta) \in W$ is allowable then for every $\beta < \delta$, $\forceP_{\beta} \Vdash| \forceP(\beta)|= \aleph_1$, thus every factor of $\forceP$ is forced to have size $\aleph_1$.
\item Every allowable forcing over $W$ is ccc and thus preserves cardinals.
\item Every allowable forcing over $W$ preserves $\CH$. Furthermore, if $\forceP= (\forceP(\alpha) \, : \, \alpha < \omega_1) \in W$ is an $\omega_1$-length iteration such that each initial segment of the iteration is allowable over $W$, then $W^{\forceP} \models \CH$.
\item The product of two allowable forcings $\forceP$ and $\forceQ$ can be densely embedded into an allowable forcing provided that $C^{\forceP} \cap C^{\forceQ}=\emptyset$.
\end{enumerate}
\end{lemma}

\begin{lemma}\label{productallowable}
Let $\alpha$ be an ordinal, assume that $W'$ is some $\alpha$-allowable generic extension of $W$, and that $\forceP^1=(\forceP^1_{\beta} \,: \,\beta < \delta_1)$ and $\forceP^2=(\forceP^2_{\beta} \,: \, \beta < \delta_2) $ are two $\alpha$-allowable forcings over $W'$ with respect to a common set $E=E_0 \cup E_1=\{(\dot{y}_{\delta},\dot{m}_{\delta}, \dot{k}_{\delta}) \,: \, \delta < \alpha\} \cup \{ (\dot{x}_{\delta},i_{\delta}) \, : \, \delta < \alpha \}$  and bookkeeping functions $F_1$ and $F_2$ respectively. Further assume that 
\[ C^{\forceP^1} \cap C^{\forceP^2} = \emptyset, \]
and that $E_1= \{ (\dot{x_{\delta}},i_{\delta}) \mid \delta < \alpha \}$ consists of reals which are in fact elements of $W'$.
Then there is a bookkeeping function $F$ such that $\forceP^1 \times \forceP^2$ densely embeds into an $\alpha$-allowable forcing over $W'$ with respect to $E$ and $F$.
\end{lemma}

\begin{proof}
We shall show that $\forceP^1 \ast \check{\forceP}^2$ is $\alpha$-allowable with respect to $E$ and some $F$. This iteration should then serve as the asserted $\alpha$-allowable forcing in which $\forceP^1 \times \forceP^2$ densely embeds. 
We define $F \upharpoonright \delta_1$ to be $F_1$. For values $\delta_1+ \beta  \ge \delta_1$ we let $F(\delta_1+\beta)$ be $F_2(\beta)$ where we identify $\forceP^2_{\beta}$-names with the $\forceP^1$-check names for $\forceP^2_{\beta}$-names. The lemma is proved by induction on $\alpha$. If $\alpha=0$, then this is Lemma \ref{FirstPropertiesOfAllowableForcings}. 

So we assume that $\alpha=1$. We assume first that, we have enlarged $E_1$ in order to get the notion of 1-allowable. Then there is a real $x_0 \in W'$ and $i \in \{0,1\}$ such that 1-allowable with respect to $E$ and some $F$ is just an allowable forcing with the additional constraint, that we must not use forcings of the form $\operatorname{Code} (x_0,y,a,b,i, \eta)$ for every $\eta$. This means that
    \begin{align*}
        \forall \eta < \delta_1 ( \forall \zeta < \omega_1 (( \mathbb{P}^1(\eta)^{G^1_{\eta}} \neq \operatorname{Code}(x_0, a, b, i, \zeta) )).
    \end{align*}
    and
    \begin{align*}
        \forall \eta < \delta_2 ( \forall \zeta < \omega_1 (( \mathbb{P}^2(\eta)^{G^2_{\eta}} \neq \operatorname{Code}(x_0, a, b, i, \zeta) )).
    \end{align*}
    This property is clearly preserved under $\forceP^1 \times \forceP^2$.
    So if both $\forceP^1$ and $\forceP^2_{\beta}$ are 1-allowable with respect to $E $ and $F_1$ and $F_2$ respectively, then, at stage $\beta$ $\forceP^1 \times \forceP^2$ is 1-allowable with respect to $F$ and $E$ as asserted.


So we can assume that 1-allowable was obtained via enlarging $E=E_0= \{(\dot{y},\dot{m},\dot{k})\}$. 
We assume first that $W= W'$. It will become clear once the proof is finished for this case that extending the proof to a $W'$ is almost trivial.
Let $(\forceP_2)_{\beta}$ be the iteration of $\forceP_2$ up to stage $\beta < \delta_2$. Assume, that $\forceP_1 \times (\forceP_2)_{\beta}$ is in fact an $\alpha$-allowable forcing relative to $E$ and $F$. Let $(G^1 \times G^2_{\beta})$ be a $\forceP^1 \times \forceP^2_{\beta}$-generic filter over $W$.  Then we have that $F(\delta_1+\beta) =F_2(\beta)  =(\dot{x},\dot{y},\dot{m},\dot{k} ,i,\dot{\eta})$, and we claim that

\begin{claim}
At any stage $\beta < \delta_2$, if we evaluate the names in $F_2(\beta)$ using $G^2_\beta$, the specific case (Case 1.1, 1.2, 1.3, 1.4 from the definition of $\alpha$-allowability) that applies in the model $W[G^2_\beta]$ is exactly the same case that applies in the extended model $W[G^1][G^2_\beta]$.
\end{claim}


If the claim holds, the lemma follows immediately by induction on $\beta < \delta_2$. Working in $W[G^1][G^2_\beta]$, we evaluate $F(\delta_1 + \beta) = F_2(\beta)$. By the claim, the case determining $\mathbb{P}^2(\beta)^{G^2_\beta}$ over the larger model $W[G^1][G^2_\beta]$ is identical to the one over $W[G^2_\beta]$, ensuring the rules of $\alpha$-allowability are strictly obeyed.



\medskip

\begin{comment}
    To prove the claim  we need to show that  $(\forceP^2_{\beta}\, : \, \beta < \delta_2)$ is still 1-allowable with respect to $E$ and $F$, when considered in the universe $W[\forceP^1]$. We fix a $\prod_{i < \omega_1} S_{i} \times \prod_{i < \omega_1} \mathbb{C} (\omega_1)$-generic filter $G \times H$ over $L$, where $G$ should be generic for the Suslin tree forcing and $H$ is generic for the $\omega_1$-Cohen forcing. Then $W$ is assumed to be a 0-allowable extension of $L[G\times H]$, and we let $\forceP$ be such that $L[G \times H] [\forceP]=W$. We note that the missing $\aleph_1$-many $h_j$'s from $H$ which are not used in the definition of $\forceP$ will not add reals to $L[G] [\{C_i \mid i \in \alpha < \omega_1 \}] [\forceP]$ by Easton's Lemma applied over $L[\{C_i \mid i \in \alpha < \omega_1 \}]$, that is 
\[ (P(\omega))^{L[G][\{C_i \mid i \in \alpha < \omega_1 \}] [\forceP]}= (P(\omega))^{L[G][H][\forceP]}.\]
So $\forceP$ can correctly be defined in an inner model of $L[G \times H]$ which only uses countably many coordinates of $H$, so $\forceP \in L[G] [ \{C_i \mid i \in \alpha < \omega_1 \}]$.
With this is mind we can assume without loss of generality that $W=W$. The proof will carry over to the case $W\ne W$ in a straightforward way.
\end{comment}


We assume first that at stage $\beta$, in the definition of $ \forceP^2$,  case 1.1 in the definition of 1-allowable applies, when working in the model $W[G^2_{\beta}]$ relative to $E$ and $F_2$. So we assume that $\beta$ is a stage such that
\[ F_2(\beta) = (\dot{x},\dot{y},\dot{m},\dot{k},i,\dot{\eta}) \]
and $(\dot{y},\dot{m},\dot{k}) \in E$ and if $\dot{x}^{G^2_{\beta}}=x$ and $\dot{y}^{G^2_{\beta}}=y$, the universe $W[ G^2_{\beta}]$ thinks that
\begin{align*}
\exists \forceQ (&\forceQ \text{ is } 0 \text{-allowable with respect to } E \text{ and some F' }\,  \land \\& \, \forceQ \Vdash {x} \in A_m({y})).
\end{align*}
Now when working over $W[G^1 \times G^2_{\beta}]$ instead, we shall show that
we are still in case 1.1.

The idea is that $ \forceQ$, as it will force $x \in A_m(y)$, will witness that we are in case 1 as well when working over $W[G^1][G^2_{\beta}]$. For that we need to show that $\forceP^1 \times \forceQ$ is 0-allowable, or in other words, that their coding areas do not intersect. As there is no reason why this should be the case, we need an additional argument. For that we show that the existence of a 0-allowable forcing which forces some $\Sigma^1_3$-statement is not tied to the coding areas it uses, as they are generically added sets. This will imply that we can shift the coding areas of $\forceQ$ away from $C^{G^1}$ and still obtain allowable forcings using these shifted coding areas instead while still forcing the $\Sigma^1_3$-assertion to hold.

We shall explain this shifting mechanism now. We write $W= L[G]$ and let $(G)_0$ and $(G)_1$ denote its projection to the first and the second coordinate respectively, so that $(G)_0$ is $\forceP^0= \prod_{i < \omega_1} \vec{S}$-generic and $(G)_1$ is $\forceP^1= \prod_{i < \omega_1} \mathbb{C}(\omega_1)^L$-generic. We write $(G)_1$ as $(C_i \mid i < \omega_1)$.  As $ G^2_{\beta}$ picks at each stage exactly one coding area $C_i$ there is $\eta < \omega_1$ and we can write $L[(G)_0][\{C_i \mid i < \omega_1 \}][G^2_{\beta}]$ as $L[(G)_0][\{C_i \mid i < \eta \}] [G^2_{\beta}] [\{C_i \mid \eta \le i < \omega_1 \}]$ and  the coding areas of $\forceP^2_{\beta}$ relative to $G^2_{\beta}$ is a subset of $\eta$.
\begin{align*}
L[(G)_0] [\{C_i \mid i < \eta \}] [G^2_{\beta}] [\{C_{\eta+ j} \mid j < \omega_1\} ] \models  \exists \forceQ (&\forceQ \text{ is } 0 \text{-allowable with} \\&
\text{respect to } E \text{ and some F' }  \\& \land \forceQ \Vdash {x} \in A_m({y})).
\end{align*}
Thus there is a condition $p \in \prod_{j<\omega_1} C_{\eta+j} \subset (\prod_{\eta \le j <\omega_1} \mathbb{C} (\omega_1))^L$
such that 
\begin{align*}
L[G] [\{C_i \mid i < \eta \}] [G^2_{\beta}] \models p  \Vdash_{(\prod_{j \in [\eta,\omega_1)} \mathbb{C} (\omega_1))^L}  \exists \forceQ (&\forceQ \text{ is } 0 \text{-allowable with} \\&
\text{respect to } E \text{ and some F' }  \\& \land \forceQ \Vdash {x} \in A_m({y})).
\end{align*}
As $(\prod_{j \in [\eta,\omega_1)} \mathbb{C} (\omega_1))^L$ is isomorphic to any 
$(\prod_{ i \in [\zeta,\omega_1)} \mathbb{C} (\omega_1))^L$, for $\zeta < \omega_1$, we can find a sufficiently large $\zeta< \omega _1$ such that the interval $[\zeta, \omega_1)$ is fully disjoint from the coding areas from $\forceP^1$ relative to $G^1$. We can also find a condition $p' \in (\prod_{j \in [\zeta,\omega_1)} \mathbb{C} (\omega_1))^L$ such that 
\begin{align*}
L[G] [\{C_i \mid i < \eta \}] [G^2_{\beta}] \models p'  \Vdash_{(\prod_{j \in [\zeta,\omega_1)} \mathbb{C} (\omega_1))^L}  \exists \forceQ (&\forceQ \text{ is } 0 \text{-allowable with} \\&
\text{respect to } E \text{ and some $F'$ }  \\& \land \forceQ \Vdash {x} \in A_m({y})).
\end{align*}
Note that $(\prod_{i< \omega_1} \mathbb{C} (\omega_1))^L$ is weakly homogeneous \footnote{The homogeneity of the forcing is in fact not necessarily needed for this argument and its use can be replaced by a density argument. We use the homogeneity as the argument becomes neater this way.}, so in fact
\begin{align*}
L[G] [\{C_i \mid i < \eta \}] [G^2_{\beta}] \models 1  \Vdash_{(\prod_{j \in [\zeta,\omega_1)} \mathbb{C} (\omega_1))^L}  \exists \forceQ (&\forceQ \text{ is } 0 \text{-allowable with} \\&
\text{respect to } E \text{ and some $F'$}  \\& \land \forceQ \Vdash {x} \in A_m({y})).
\end{align*}
So for any $\zeta < \omega_1$, $\zeta > \eta$ we can find a condition $q \in (\prod_{j \in [\zeta,\omega_1) } \mathbb{C} (\omega_1))^L$ such that $q \in \prod_{j \in [\zeta,\omega_1) } C_j$ and
\begin{align*}
L[G] [\{C_i \mid i < \eta \}] [G^2_{\beta}] \models q  \Vdash_{(\prod_{j \in [\zeta,\omega_1)} \mathbb{C} (\omega_1))^L}  \exists \forceQ (&\forceQ \text{ is } 0 \text{-allowable with} \\&
\text{respect to } E \text{ and some $F'$ }  \\& \land \forceQ \Vdash {x} \in A_m({y})).
\end{align*}
As the coding areas of $\forceP^1$ are countable, we can find a $\zeta> \eta$ such that $C^{\forceP^1} \cap [\zeta, \omega_1) = \emptyset$. Also the coding areas of $\forceQ$ are contained as a subset in the interval $[\zeta,\omega_1)$ so we can conclude
\begin{align*}
L[(G)_0] [(G)_1] [G^2_{\beta}] \models  \exists \forceQ (&\forceQ \text{ is } 0 \text{-allowable with} \\&
\text{respect to } E \text{ and some $F'$ }  \\&  \land C^{\forceQ} \cap C^{\forceP^1}=\emptyset \\& \land
\forceQ \Vdash {x} \in A_m({y})).
\end{align*}
and so
\begin{align*}
L[(G)_0] [(G)_1] [G^1] [G^2_{\beta}] \models  \exists \forceQ (&\forceQ \text{ is } 0 \text{-allowable with} \\&
\text{respect to } E \text{ and some $F'$ }  \\&  \land
\forceQ \Vdash {x} \in A_m({y})).
\end{align*}


%But $p'$ is not necessarily in our fixed $H= \prod_{i < \omega_1} C_i$. So we need to argue a bit more.



%For $p \in \prod_{i \in [\alpha, \omega_1)} \mathbb{C} (\omega_1) \subset \prod_{\omega_1} \mathbb{C} (\omega_1)$
%there is always a dense set $D_p:= \{ q \in \prod_{i < \omega_1} \mathbb{C} (\omega_1) \mid \exists \{\alpha_0,...,\alpha_{\omega},... \} (\{ \alpha_0...\} \cap \text{dom} (p) = \emptyset \land \{\alpha_0,...\} \subset \text{dom} (q) \land q \upharpoonright \{ \alpha_0,...,\} \cong p \}.$

%So in $L[G] [\{C_i \mid i < \alpha \} ] [\forceP^2_{\beta}]$,  $q \upharpoonright \{ \alpha_0,...\} \cong p$ and $q \upharpoonright \{ \alpha_0 ...\} \Vdash \exists \forceQ (\forceQ \text{ is } 0 \text{-allowable with}
%\text{ respect to } E \text{ and some F } \land \forceQ \Vdash {x} \in A_m({y})), $ and moreover $\forceQ$ only uses elements from $\{\alpha_0,...,\}$ as coding area.
%So also as $q < q \upharpoonright \{ \alpha_0...\}$ we know that also
%in $L[G] [\{C_i \mid i < \alpha \} ] [\forceP^2_{\beta}]$
%$q \Vdash \exists \forceQ (\forceQ $ is 0-allowable with
 %respect to $ E$  and some $F$ and $ \forceQ \Vdash {x} \in A_m({y}))$ and $\forceQ$ %uses only coding areas from $\{\alpha_0,...\}$.
 
 

%As $p \in H$ there is a $q \in D_p \cap H$. Let $r < q, p$ such that $r \in H$. Then
%\begin{align*}
%r \Vdash_{\prod_{\omega_1} \mathbb{C} (\omega_1)} \exists \forceQ \text{ with coding areas contained in $\{\alpha_0...\}$ and } \forceQ \Vdash x \in A_m(y) 
%\end{align*}

%So, as $r \in H$
%\begin{align*}
%L[G ][ H] [\forceP^2_{\beta} ]\models \exists \forceQ \text{ with coding areas %contained in $\{\alpha_0...\}$ and } \forceQ \Vdash x \in A_m(y) 
%\end{align*}
%However $\{\alpha_0,...\}$ is disjoint from the coding areas from $\forceP^1$, which is what we wanted. So $\forceP^1 \times \forceQ$ is 0-allowable and forces $x \in A_m(y)$.

So at stage $\beta$, we are in case 1 as well, when working over $W[G^1][G^2_{\beta}]$ which proves the Claim in the first instance.

If, at stage $\beta$ and working in the model $W[G_{\beta}^2]$ case 2 applies, when considering $ \forceP_2$ as a 1-allowable forcing with respect to $E$ and $F^2$ over $W$, then we argue first that case 1 is impossible when considering $\forceP_2$ as a $1$-allowable forcing over $W[G^1]$, where as $G^1$ is $\forceP^1$-generic as before.


Indeed, assume for a contradiction that case 1 must be applied at stage $\beta$, then, by definition,
\begin{align*}
W [G^1] [G^2_{\beta}] \models  \exists \forceQ (&\forceQ \text{ is } 0 \text{-allowable with} \\&
\text{respect to } E \text{ and some $F'$}  \\&  \land
\forceQ \Vdash {x} \in A_m({y})).
\end{align*}
But then $\forceP^1 \times \forceQ$ is such that there is a condition $p^1 \in \forceP^1 \cap G^1$ which forces that $p^1 \Vdash C^{\forceP^1} \cap C^{\forceQ} = \emptyset$. So restricting $\forceP^1 \times \forceQ$ to conditions stronger than $(p^1,1)$ will yield a $0$-allowable forcing $(\forceP^1 \times \forceQ)_{\le (p^1,1)} $ such that 
\begin{align*}
W [G^2_{\beta}] \models  (&(\forceP^1 \times \forceQ)_{\le (p^1,1)} \text{ is } 0 \text{-allowable with} \\&
\text{respect to } E \text{ and some $F'$ }  \\&  \land
(\forceP^1 \times \forceQ)_{\le (p^1,1)} \Vdash {x} \in A_m({y})).
\end{align*}
So we are in case 1 as well when working over $W [G^2_{\beta}]$, which is a contradiction to our assumption.

Now we need to argue that we are in case 2, when working over $W[G^1][G^2_{\beta}]$ in the definition of $\forceP^2$. We can argue in a very similar way than before. So suppose we are at stage $\beta$ of our iteration and we work over the model $W[G^1][G^2_{\beta}]$ and case 1 does not apply, which we know already. The universe $W[G^2_{\beta}]$ thinks that there is $\forceQ$ which is 0-allowable and which forces $x \in A_k(y)$. Just as before we would like to form $\forceP^1 \times \forceQ$ in the model $W[G^2_{\beta}]$ which will force $x \in A_k(y)$, therefore witnessing that we are in case 2, which is what we want to finish the proof. Again, as before, there is no guarantee that $\forceP^1$ and $\forceQ$ have non-intersecting coding areas, so we need to apply again the shifting argument from above.
As there we can argue that there is in fact an allowable $\forceQ$ such that $C^{\forceQ} \cap C^{\forceP^1}= \emptyset$ and such that $\forceQ \Vdash x \in A_k (y)$. Now this $\forceQ$ can be taken to form a product with $\forceP^1$, which witnesses that we are in case 2 when defining $\forceP^2$ at stage $\beta$ over the model $W[G^2_{\beta}]$ so this proves the claim for case 2 in the definition of allowable forcings.

Finally, if at stage $\beta$, case 3 applies when considering $\forceP_2$ as a 1-allowable forcing with respect to $E$ over $W[G^2_{\beta}]$, we claim that we must be in case 3 as well, when defining $\forceP_2 (\beta)$ over $W[G^1][G^2_{\beta}]$.

If not, then we would be in case 1 or 2 when defining $\forceP^2(
\beta)$ over $W[G^1][G^2_{\beta}]$. Assume without loss of generality that we were in case 1, then there would be a 0-allowable $\forceQ \in W[G^1][G^2_{\beta}]$ such that $\forceQ \Vdash x \in A_m(y)$. Again, via shifting the coding areas of $\forceQ$ we can assume that the forcing $\forceP^1 \ast\forceQ $ in $W[G^2_{\beta}]$ is allowable and witnesses that we are in fact in case 1 as well when defining  $\forceP^2(\beta)$ over $W[G^2_{\beta}]$ which is a contradiction to our assumption that we are in case 3. This shows that we must be in case 3 as well, and the claim, and hence the lemma is finally proved for $\alpha=1$.

 

\medskip

We shall argue now that the Claim is true for $\alpha+1$-allowable forcings provided we know that it is true for $\alpha$-allowable forcings. Again we can focus on the case when $\alpha+1$ allowable forcings are obtained via enlarging $E_0$, as enlarging $E_1$ just means to avoid certain coding forcings, which is trivial to be closed under products.
We shall show the claim via induction on $\beta$. The proof itself is almost identical to the case $\alpha=1$. As there, we can shift around coding areas using isomorphisms of $(\prod_{\omega_1} \mathbb{C} (\omega_1))^L$ to always find the relevant $\alpha$-allowable forcings whose coding areas are disjoint. All the work we have to do is replace $0$-allowable with $\alpha$-allowable.



This ends the proof of the claim and so we have shown the lemma.
\end{proof}

\end{document}
